% Chapter 1. Update the research questions or hypotheses only if the study design changes.
\chapter{Introduction}
\label{ch:intro}

\section{Motivation}

Nitrogen dioxide (\ce{NO2}) forms mainly in combustion, for example in vehicle engines and power
plants. It irritates the airways, and long exposure is linked to respiratory disease. In 2021 the
World Health Organization lowered its guideline for the annual mean concentration of \ce{NO2} to
10~\(\mu\)g\,m\(^{-3}\) and added a 24-hour guideline of 25~\(\mu\)g\,m\(^{-3}\)
\cite{who2021aqg}. Checking such low levels across a city needs many sensors, so the sensors must
be cheap, small and able to work at room temperature with little power. A common design is the
chemiresistor: a thin sensing layer whose electrical resistance changes when gas molecules adsorb
on it.

Graphene is an attractive sensing layer because every atom is on the surface. Schedin et
al.\ showed that a graphene device can detect the adsorption of single \ce{NO2} molecules
\cite{schedin2007detection}. \ce{NO2} takes electrons from graphene, which changes the number of
charge carriers and therefore the resistance. Pristine graphene binds most gases weakly and has
no band gap, so its response and its selectivity are limited \cite{leenaerts2008adsorption}.
Changing the chemistry of the carbon layer, for example with dopants or defects, changes how
strongly each gas binds \cite{zhang2009dopants,yuan2013graphene}.

Graphene quantum dots (GQDs) are flakes of graphene a few nanometres across, with hydrogen atoms
on the edge carbons. Their finite size opens a gap between the highest occupied and the lowest
unoccupied molecular orbitals (HOMO and LUMO), and the size of this gap depends on the size,
shape and edge type of the flake \cite{ponomarenko2008,bacon2014gqd}. Bottom-up organic synthesis
can now make nanographenes with an exactly defined structure \cite{yan2010large,narita2015nanographene},
so it makes sense to ask which exact structure is best for a given purpose. GQDs have already
been used in \ce{NO2} sensors; for example, GQDs on \ce{SnO2} nanodomes gave a detection limit of
1.1~ppb at room temperature \cite{lee2023gqdsno2}. What is missing is a systematic way to choose
the GQD structure itself.

\section{The design problem}

The number of possible GQDs grows quickly with the number of fused hexagons. The dataset used
in this work contains 91{,}174 distinct GQDs with 20 to 40 hexagons, and it is only a sample of
what is possible at these sizes. A useful sensor material has to meet three requirements at once.
It must respond to \ce{NO2}, which means that \ce{NO2} should take a noticeable amount of charge
from it. It must recover, which means that \ce{NO2} must not bind so strongly that it stays on the
surface. And it must be selective, which means that \ce{NO2} must bind more strongly than the other
gases in air. These requirements pull in different directions: a GQD that gives electrons easily
also tends to form a chemical bond with \ce{NO2}, and then the sensor does not reset.

Each candidate has to be evaluated with quantum chemistry. In this work, one evaluation with the
semiempirical method GFN2-xTB \cite{bannwarth2019gfn2} takes about a minute on a laptop, and a
check with density functional theory (DFT) takes about an hour per GQD on a data-centre GPU. With
this cost, exhaustive screening is out of reach, and the search method has to find good structures
with a small number of evaluations. In molecular design this property is called sample efficiency
\cite{gao2022pmo}.

\section{Approach}

Large language models (LLMs) can read and reason about chemistry, but they are poor at writing
long molecular strings such as SMILES correctly \cite{guo2023llmchem}. Recent work has used LLMs as
the mutation operator inside evolutionary algorithms, where the LLM proposes changes to good
candidates and an objective function decides which changes survive
\cite{lehman2022elm,romeraparedes2024funsearch,wang2024molleo}. Other work has given LLMs access
to chemistry software as tools, so that the model plans while the tools compute
\cite{yao2023react,bran2024chemcrow,boiko2023coscientist}.

This thesis combines these ideas with a strict division of work: \emph{the LLM suggests, and code
and quantum chemistry decide}. The LLM describes a GQD as a list of hexagon positions on a
honeycomb grid instead of a SMILES string. Code turns this list into a molecule, checks it, and
explains any problem in terms the LLM can act on. A score computed with GFN2-xTB measures how good
the GQD is as an \ce{NO2} sensor, and an evolutionary loop asks the LLM to edit the best GQDs found
so far. The same loop with random edits is the control. DFT calculations check where the
xTB numbers can be trusted.

The scope is limited to pure-carbon, closed-shell GQDs with 20 to 40 hexagons and no holes. All
calculations describe a single gas molecule on a single GQD in the gas phase. Dopants, edge groups,
fine-tuning of the LLM and dedicated generative models are left for future work
(\cref{ch:conclusion}).

\section{Research questions}
\label{sec:rq}

\begin{enumerate}[label=\textbf{RQ\arabic*.},leftmargin=*]
  \item Can a general-purpose LLM design valid GQDs when it writes hexagon positions instead of
        SMILES, and what prompt and feedback make it reliable?
  \item Can a score computed with GFN2-xTB rank GQDs for \ce{NO2} sensing, and how well do its
        numbers agree with DFT?
  \item With the same number of xTB evaluations, does an LLM used as the mutation operator find
        better \ce{NO2}-sensing GQDs than random mutation?
  \item What do the best GQDs have in common?
\end{enumerate}

\section{Hypotheses}
\label{sec:hypotheses}

\begin{enumerate}[label=\textbf{H\arabic*.},leftmargin=*]
  \item Most invalid LLM designs fail because they are open-shell, a global property that is hard
        to see from coordinates. If the prompt states the closed-shell rule and the feedback lists
        concrete edits that fix a rejected design, more than 90\% of requests will end in a valid
        design within two retries. (Tested in experiment E1, \cref{sec:e1}.)
  \item xTB will agree with DFT for physisorbed gases and will rank GQDs by charge transfer in
        the same order as DFT, even where the absolute values differ. It will be less reliable for
        chemisorbed \ce{NO2} and for \ce{O2}. Across a benchmark of GQDs, we expect a Spearman rank
        correlation of at least 0.7 between the xTB and DFT values of the sensing descriptors.
        (Experiments E2 to E4.)
  \item The LLM sees the structures and scores of the best GQDs and the outcome of its earlier
        edits, so its edits will be better than chance. Across repeated runs with the same xTB
        budget, the LLM loop will reach a higher best score and a higher mean score of the
        evaluated children than random mutation. (Experiment E5.)
  \item A smaller HOMO-LUMO gap raises the charge given to \ce{NO2} but also the chance that
        \ce{NO2} forms a covalent bond, which blocks recovery. The best sensors will therefore have
        an intermediate gap and physisorb \ce{NO2}. (Experiments E2 and E5.)
\end{enumerate}

\section{Contributions}

The work makes six contributions.
\begin{enumerate}
  \item A hexagon-grid representation of GQDs with exact conversion between hexagon lists and
        molecules. The conversion checks the closed-shell condition directly and, for an
        open-shell design, lists the hexagon edits that fix it. All 91{,}174 dataset structures
        convert and convert back without error.
  \item A machine-learned predictor of the GFN2-xTB gap (mean absolute error 0.023~eV on held-out
        GQDs) with a calibrated uncertainty, used as a free pre-filter.
  \item An \ce{NO2} sensing score from GFN2-xTB, calibrated on 50 GQDs and checked against DFT.
  \item An evolutionary loop with an LLM or random mutation operator, code-based validation and
        repair, a diversity rule, and a log of every proposal.
  \item Controlled experiments: an ablation of prompt and feedback for LLM design, a comparison of
        LLM and random mutation under the same xTB budget, and a GPU-accelerated DFT benchmark.
  \item A log of edits and their measured effects, which is the starting point for fine-tuning a
        smaller model.
\end{enumerate}

\section{Thesis outline}

\Cref{ch:literature} reviews GQDs, gas sensing with graphene, the computational methods used here,
and earlier work on molecular design with machine learning and LLMs. \Cref{ch:methods} describes
the workflow and its algorithms in the order they were built. \Cref{ch:experiments} defines the
experiments that test each hypothesis. \Cref{ch:results} reports the results, \cref{ch:discussion}
discusses what they mean and where they are limited, and \cref{ch:conclusion} concludes and lists
the next steps. The appendices give the LLM prompts, the logged fields and the commands needed to
reproduce the work.
